% GENERATED FILE. Edit canonical lessons, then run npm run generate:books.
% canonical-source-hash: fnv1a64:e359b96c10b0f4f8
% canonical-lessons: MW-C132-achar, MW-C132-lapsi, MW-C132-ghevar, MW-C132-makki, MW-C132-gehu

\chapter{Pickle, Sweet wheat dish, Honeycomb sweet, Maize, Wheat}
\label{ch:mw-pickle-sweet-wheat-dish-honeycomb-sweet-maize-wheat}
\hlchaptermodality{132}

\begin{chapteropening}
\textbf{By the end of this chapter:} \emph{I can say pickle, a sweet wheat dish, a honeycomb sweet, maize and wheat.}

Complete the last lesson of chapter 132: i can say pickle, a sweet wheat dish, a honeycomb sweet, maize and wheat.
\end{chapteropening}

\section[\texorpdfstring{\mw{अचार} (achār)}{अचार (achār)}]{\mw{अचार} (achār) — pickle}

\label{lesson:MW-C132-achar}



\begin{quote}
Before the new one: say the Marwadi for gram-flour dumplings, then the Marwadi for a papad.
\end{quote}

\subsection*{You'll want to know: \mw{अचार}}
\textbf{\mw{अचार}} — \emph{achār} — \textquotedblleft{}pickle\textquotedblright{}.

\begin{grammarlens}[title={What you've built}]
One more word for this chapter.
\end{grammarlens}

\subsection*{Your turn}
\begin{itemize}
\raggedright
  \item \emph{Say it:} \emph{achār}
  \item \emph{Say it:} \emph{achār}, once more
  \item \emph{Say it:} say \emph{achār}
  \item \emph{Recall:} say the Marwadi for gram-flour dumplings, then the Marwadi for a papad, then say \emph{achār} again
\end{itemize}

\begin{tcolorbox}[breakable,colback=teal!4,colframe=teal!35!black,title={Before you move on}]
What does \mw{अचार} mean? (\textquotedblleft{}Pickle\textquotedblright{}.) Say it once more.
\end{tcolorbox}
\section[\texorpdfstring{\mw{लापसी} (lāpsī)}{लापसी (lāpsī)}]{\mw{लापसी} (lāpsī) — a sweet wheat dish}

\label{lesson:MW-C132-lapsi}



\begin{quote}
Before the new one: say the Marwadi for a papad, then the Marwadi for pickle.
\end{quote}

\subsection*{You'll want to know: \mw{लापसी}}
\textbf{\mw{लापसी}} — \emph{lāpsī} — \textquotedblleft{}a sweet wheat dish\textquotedblright{}.

\begin{grammarlens}[title={What you've built}]
One more word for this chapter.
\end{grammarlens}

\subsection*{Your turn}
\begin{itemize}
\raggedright
  \item \emph{Say it:} \emph{lāpsī}
  \item \emph{Say it:} \emph{lāpsī}, once more
  \item \emph{Say it:} say \emph{lāpsī}
  \item \emph{Recall:} say the Marwadi for a papad, then the Marwadi for pickle, then say \emph{lāpsī} again
\end{itemize}

\begin{tcolorbox}[breakable,colback=teal!4,colframe=teal!35!black,title={Before you move on}]
What does \mw{लापसी} mean? (\textquotedblleft{}A sweet wheat dish\textquotedblright{}.) Say it once more.
\end{tcolorbox}
\section[\texorpdfstring{\mw{घेवर} (ghevar)}{घेवर (ghevar)}]{\mw{घेवर} (ghevar) — a honeycomb sweet}

\label{lesson:MW-C132-ghevar}



\begin{quote}
Before the new one: say the Marwadi for pickle, then the Marwadi for a sweet wheat dish.
\end{quote}

\subsection*{You'll want to know: \mw{घेवर}}
\textbf{\mw{घेवर}} — \emph{ghevar} — \textquotedblleft{}a honeycomb sweet\textquotedblright{}.

\begin{grammarlens}[title={What you've built}]
One more word for this chapter.
\end{grammarlens}

\subsection*{Your turn}
\begin{itemize}
\raggedright
  \item \emph{Say it:} \emph{ghevar}
  \item \emph{Say it:} \emph{ghevar}, once more
  \item \emph{Say it:} say \emph{ghevar}
  \item \emph{Recall:} say the Marwadi for pickle, then the Marwadi for a sweet wheat dish, then say \emph{ghevar} again
\end{itemize}

\begin{tcolorbox}[breakable,colback=teal!4,colframe=teal!35!black,title={Before you move on}]
What does \mw{घेवर} mean? (\textquotedblleft{}A honeycomb sweet\textquotedblright{}.) Say it once more.
\end{tcolorbox}
\section[\texorpdfstring{\mw{मक्की} (makkī)}{मक्की (makkī)}]{\mw{मक्की} (makkī) — maize}

\label{lesson:MW-C132-makki}



\begin{quote}
Before the new one: say the Marwadi for a sweet wheat dish, then the Marwadi for a honeycomb sweet.
\end{quote}

\subsection*{You'll want to know: \mw{मक्की}}
\textbf{\mw{मक्की}} — \emph{makkī} — \textquotedblleft{}maize\textquotedblright{}.

\begin{grammarlens}[title={What you've built}]
One more word for this chapter.
\end{grammarlens}

\subsection*{Your turn}
\begin{itemize}
\raggedright
  \item \emph{Say it:} \emph{makkī}
  \item \emph{Say it:} \emph{makkī}, once more
  \item \emph{Say it:} say \emph{makkī}
  \item \emph{Recall:} say the Marwadi for a sweet wheat dish, then the Marwadi for a honeycomb sweet, then say \emph{makkī} again
\end{itemize}

\begin{tcolorbox}[breakable,colback=teal!4,colframe=teal!35!black,title={Before you move on}]
What does \mw{मक्की} mean? (\textquotedblleft{}Maize\textquotedblright{}.) Say it once more.
\end{tcolorbox}
\section[\texorpdfstring{\mw{गेहूं} (gehū̃)}{गेहूं (gehū̃)}]{\mw{गेहूं} (gehū̃) — wheat}

\label{lesson:MW-C132-gehu}



\begin{quote}
Before the new one: say the Marwadi for a honeycomb sweet, then the Marwadi for maize.
\end{quote}

\subsection*{You'll want to know: \mw{गेहूं}}
\textbf{\mw{गेहूं}} — \emph{gehū̃} — \textquotedblleft{}wheat\textquotedblright{}.

\begin{grammarlens}[title={What you've built}]
One more word for this chapter.
\end{grammarlens}

\subsection*{Your turn}
\begin{itemize}
\raggedright
  \item \emph{Say it:} \emph{gehū̃}
  \item \emph{Say it:} \emph{gehū̃}, once more
  \item \emph{Say it:} say \emph{gehū̃}
  \item \emph{Recall:} say the Marwadi for a honeycomb sweet, then the Marwadi for maize, then say \emph{gehū̃} again
\end{itemize}

\begin{tcolorbox}[breakable,colback=teal!4,colframe=teal!35!black,title={Before you move on}]
What does \mw{गेहूं} mean? (\textquotedblleft{}Wheat\textquotedblright{}.) Say it once more.
\end{tcolorbox}
